%%%%%%%%%%%%%%%%%%%%%%%%%%%%%%%%%
% Acknowledgements              %
%%%%%%%%%%%%%%%%%%%%%%%%%%%%%%%%%
%
% This page contains the acknowledgements section of your report. The acknowledgements section is used to express gratitude and appreciation to individuals, institutions, or organizations that have provided support, guidance, or assistance during the completion of your project.

%%%%%%%%%%%%%%%%%%%%%%%%%%%%%%%%%

% \centerline{
% Uncomment the line above to center the section heading

\section*{Acknowledgements}

% }
% Uncomment the closing bracket above to center the section heading

\addcontentsline{toc}{section}{Acknowledgements}
% Add the Acknowledgements section to the Table of Contents

\vspace{1.4cm}
\noindent
% Start a new paragraph without indentation

The project team wishes to express sincere gratitude to all those who made this work possible.

We are deeply grateful to our academic supervisor, \textbf{Dr.\ Mortadha Dahmani}, for his continuous guidance, technical advice, and institutional support throughout the project. His provision of VPS infrastructure and dedicated office hours were instrumental in enabling the cloud connectivity features of this system.

We thank \textbf{BAKO Motors} and our institution supervisor, \textbf{Achraf}, for granting access to their electric motorcycle prototype and for sharing domain knowledge about the vehicle's CAN bus architecture. The weekly on-vehicle testing sessions held on the MedTech campus were critical for validating the monitoring system under real operating conditions.

We acknowledge the \textbf{MedTech Mediterranean Institute of Technology FABLAB} for providing unrestricted access to soldering stations, bench power supplies, logic analysers, and PCB rework equipment throughout the project lifecycle. This support significantly reduced prototype development costs and timelines.

Finally, the team thanks all peers and faculty members who provided feedback during progress reviews, and the open-source communities behind the ESP32 Arduino framework, KiCad, FastAPI, and the many libraries that form the backbone of this system.

